\section{Introduction}
\label{a:sec:intro}

The quasi-stationary theory of \ChM reduces the late-time conditional behaviour of the
subcritical binary Galton--Watson process to a single constant
\begin{equation}
\label{a:eq:Apdef}
A(p) = \lim_{n\to\infty}\frac{S_n}{(2p)^n},
\qquad p\in[0,\tfrac12),
\end{equation}
whose reciprocal is the limiting conditional mean. Approximate values are easy enough
to obtain by iterating the survival recursion, but iteration fails as
$p\to\tfrac12^-$ because the number of generations needed for near-stationarity grows
without bound. No formula in elementary functions is known for $A(p)$.

This chapter develops the analytic theory of that constant: an exact infinite product;
an exact series with two-sided bounds; the near-critical asymptotic; the comparison with
the continuous-time constant $\Ac(p)$; a record of closed-form searches; the
identification of $A(p)$ with a value of the Koenigs linearising function of the logistic
map; and a careful account of what hypertranscendence does and does not imply about
$A(p)$, together with inverse-symbolic evidence. Standing setup is recalled briefly from
\ChM; proofs of Galton--Watson extinction and of the conditional mean and variance are
not repeated.

The results of \cref{a:sec:product,a:sec:Abounds,a:sec:asymptotic,a:sec:koenigs} are
original to this thesis in the form presented here. Hypertranscendence claims for the
conjugacy follow Becker--Bergweiler \cite{Becker1995}, stated and applied with the scope
restrictions of \cref{a:sec:scope}.
